\section{Banc d'essai et comparaison avec les autres approches}\label{sec-banc}

\begin{objectif}
Mesurer, sur un marché Totem synthétique construit contributeur par contributeur, ce que le détecteur structurel
apporte par rapport à des contrôles naïfs et aux idées de départ : B-spline et mélange gaussien, petits réseaux de
neurones.
\end{objectif}

\subsection{Conception}
Le marché synthétique est construit à partir des contributeurs, pour que chaque mécanisme de la \cref{sec-diagnostic}
apparaisse comme dans Totem, et non comme une forme ajoutée à une courbe. Les détails sont en \cref{app-synthetic}.
\begin{itemize}
\item \textbf{Quatre catégories} (\cref{tab-categories}), onze maturités d'un mois à cinq ans, six périodes de
dividendes (la fin de 2026, puis 2027 à 2031), date de valorisation le 14~octobre 2026.
\item \textbf{Les contributeurs} ont des vues cohérentes (\cref{hyp-coherent}) : des dividendes formés d'un facteur
persistant et d'un bruit propre à chaque dividende, dont la dispersion croît avec la maturité, et une courbe de repo
lisse. Les contributeurs options valorisent calls et puts à cinq strikes avec leur propre forward et une nappe de
volatilité bruitée. Les panels se réduisent avec la maturité et sont emboîtés ; les panels options et dividendes sont
tirés indépendamment.
\item \textbf{Le consensus} est la moyenne après un filtre médiane/écart absolu médian ; dispersions et nombres de
contributeurs sont publiés comme le ferait Totem. La chaîne du contrôleur s'exécute ensuite sans modification.
\item \textbf{Les marquages FO} sont les dividendes vrais avec un bruit de \pct{1}, et le repo vrai avec une erreur de
niveau de $3$\,pb (un écart-type).
\item \textbf{Les scénarios} (\cref{tab-scenarios}) sont simulés pour chaque sous-jacent de test avec le même tirage de
contributeurs : les différences entre scénarios sont appariées, et le rapport signal sur bruit de chaque injection
est connu exactement.
\end{itemize}

\begin{table}[ht]
\centering
\caption{Les quatre catégories du marché synthétique.}\label{tab-categories}
\input{generated/\lang/tab_categories}
\end{table}

\begin{table}[ht]
\centering
\caption{Scénarios du banc d'essai. Les six premiers sont des anomalies de la courbe Totem ; pour les deux derniers,
une alerte d'anomalie Totem n'est pas le résultat souhaité (voir le texte).}\label{tab-scenarios}
\small
\begin{tabular}{@{}>{\raggedright\arraybackslash}p{3.4cm}>{\raggedright\arraybackslash}p{7.5cm}
>{\raggedright\arraybackslash}p{3.3cm}@{}}
\toprule
Scénario & Injection & Amplitude\\
\midrule
Décalage du spot & les contributeurs options valorisent avec un autre spot & $\pm\pct{0.03}$ à $\pm\pct{0.15}$\\
Erreur sur une période & une période (2027 à 2029) erronée dans le consensus dividendes & $\pm\num{0.15}$ à
$\pm\num{0.5}$ pour un spot de $100$\\
Désaccord de date & \pct{35} à \pct{70} du panel options placent de l'autre côté d'une échéance le dividende le plus
proche & un dividende\\
Profil périmé & profil d'allocation décalé pour qu'un dividende franchisse une échéance & $3$ à $20$ jours au-delà\\
Forward aberrant & calls de consensus décalés à une échéance & $3$ à $6$ écarts-types du bruit de portage\\
Sous-panel long biaisé & les contributeurs restant après la plus forte sortie sont biaisés sur les dividendes suivants
& $\pm\pct{8}$ à $\pm\pct{20}$\\
\midrule
Erreur de marquage FO & repo FO décalé parallèlement, Totem intact & $\pm15$ à $\pm50$\,pb\\
Biais brut/net & toutes les périodes Totem multipliées par $1+\tau$ & $|\tau|$ de \pct{8} à \pct{18}\\
\bottomrule
\end{tabular}
\end{table}

\paragraph{Détecteurs.} Cinq détecteurs sont comparés :
(1) \emph{écart naïf} : le plus grand $|q^{\tot}_k-q^{\fo}_k|$ ;
(2) \emph{forme naïve} : le plus grand écart du repo Totem à l'interpolation de ses deux voisins ;
(3) \emph{B-spline et mélange gaussien sur la courbe} : les six coefficients d'une B-spline cubique ajustée sur la
courbe Totem (nœuds en $\ln T$), regroupés par un mélange gaussien par catégorie dont le nombre de composantes est
choisi par le critère BIC ; le score est l'opposé de la log-vraisemblance ;
(4) \emph{la même chose sur l'écart Totem--FO} ;
(5) \emph{le détecteur structurel} de la \cref{sec-detection}.

\paragraph{Protocole.} \vNTrain{} sous-jacents sains par catégorie calibrent chaque détecteur à \vFPR{} de fausses
alertes (un quantile du score pour les quatre premiers ; $\kappa$ et l'échelle de niveau pour le détecteur
structurel). Puis \vNTestPerCat{} nouveaux sous-jacents par catégorie sont simulés sous chacun des \vNScen{}
scénarios, soit \vNTest{} sous-jacents de test.

\subsection{Résultats}
\begin{table}[ht]
\centering
\caption{Part des sous-jacents signalés, chaque détecteur étant calibré à \vFPR{} de fausses alertes sur l'historique
sain. La ligne \og Sain \fg{} est le taux de fausses alertes réalisé sur de nouveaux sous-jacents. Pour les deux
dernières lignes avant la moyenne, une alerte d'anomalie Totem n'est pas le résultat souhaité.}\label{tab-detection}
\input{generated/\lang/tab_detection}
\end{table}

\begin{figure}[ht]
\centering\includegraphics[width=\linewidth]{figures/\lang/power.pdf}
\caption{Taux de détection en fonction du rapport signal sur bruit $\lambda$ de l'anomalie injectée, sur l'ensemble des
six anomalies Totem. La courbe grise est la puissance théorique d'une forme type connue au seuil typique
(\cref{cor-power}).}
\label{fig-power}
\end{figure}

\begin{figure}[ht]
\centering\includegraphics[width=\linewidth]{figures/\lang/attribution.pdf}
\caption{Mécanisme cité en premier par le détecteur structurel, par scénario. Niveau, marche, pic et saison sont les
familles de formes types ; \og rugosité \fg{} signifie que seul le test global s'est déclenché ; \og pas d'alerte \fg{}
que le sous-jacent n'a pas été signalé.}
\label{fig-attribution}
\end{figure}

\paragraph{Lecture des résultats.}
\begin{itemize}
\item \textbf{En moyenne, les meilleures méthodes font jeu égal.} Sur les six anomalies Totem, le détecteur structurel
signale \vMeanStruct{} des sous-jacents, le mélange gaussien sur l'écart \vMeanGmmDiff, sur la courbe \vMeanGmmCurve, la
forme naïve \vMeanShape{} et l'écart naïf \vMeanNaive.
\item \textbf{Les forwards aberrants} sont le cas où le portage et le modèle de bruit rapportent le plus : \vFwdStruct{}
de détection, contre \vFwdNaive{} pour l'écart naïf, qui ne voit pas une valeur aberrante de trois à six écarts-types
sous le bruit du long terme, et \vFwdShape{} pour la forme naïve.
\item \textbf{Décalage du spot} (\vSpotStruct) et \textbf{profil périmé} (\vSeasonStruct) sont faciles pour toutes les
méthodes ; le détecteur structurel ajoute l'explication (\cref{fig-attribution}).
\item \textbf{Les erreurs sur une période} (\vBucketStruct, contre \vBucketNaive{} pour l'écart naïf) sont limitées par
leur taille : $\num{0.15}$ à $\num{0.5}$ pour un spot de $100$ est petit devant la dispersion des dividendes lointains,
et la \cref{fig-power} montre que les cas manqués ont un faible rapport signal sur bruit.
\item \textbf{Les désaccords de date} sont détectés dans \vTimingStruct{} des cas ; la forme naïve (\vTimingShape) est
compétitive, car ces pics sont grands en repo.
\item \textbf{Le sous-panel long biaisé} échappe à toutes les méthodes (\vPanelStruct{} pour le détecteur structurel,
\vPanelNaive{} pour l'écart naïf) : le biais s'accumule de façon lisse et se confond avec une pente de repo. C'est le
\cref{thm-nonid}, pas un défaut de réglage.
\item \textbf{Les vraies erreurs de marquage FO} ne sont pas prises pour des courbes Totem aberrantes : \vFoStruct{}
d'alertes (le taux sain), et \vFoMismark{} sont signalées comme erreurs de marquage FO possibles, contre \vFoClean{}
des sous-jacents sains. L'écart naïf et le mélange gaussien sur l'écart déclenchent une alerte de courbe aberrante sur
\vFoNaive{} et \vFoGmmDiff{} d'entre elles (\cref{prop-gmm-fo}). Le test d'erreur de marquage trouve \vFoMismarkEuStock{}
des erreurs sur les actions européennes et \vFoMismarkIndex{} sur l'indice, mais seulement \vFoMismarkHtb{} sur les
titres difficiles à emprunter : les vues de repo des contributeurs y diffèrent d'environ $60$\,pb, si bien qu'un écart
FO de $15$ à $50$\,pb reste dans l'incertitude du consensus lui-même.
\item \textbf{Le biais brut/net} est signalé dans \vGrossStruct{} des cas, surtout comme des marches aux détachements
des payeurs discrets ; pour le flux de dividendes d'un indice, il ressemble à une pente de repo
(\cref{prop-signatures}, (vi)).
\item \textbf{Puissance.} Pour un rapport signal sur bruit entre $4$ et $6$, le détecteur structurel trouve
\vPowStructE{} des anomalies, et au-delà de $6$ de \vPowStructF{} à \vPowStructH, conformément au \cref{cor-power}. En
dessous de $3$, l'écart naïf et le mélange gaussien sur l'écart détectent davantage (\vPowGmmDiffC{} contre
\vPowStructC{} entre $2$ et $3$) : ils réagissent aussi à la partie d'une anomalie qui ressemble à un écart de repo
lisse, que le détecteur structurel laisse volontairement au test d'erreur de marquage FO. C'est la même propriété qui
leur fait signaler les vraies erreurs de marquage FO.
\item \textbf{Explication.} Quand le détecteur structurel signale une anomalie Totem, sa première explication cite le
mécanisme injecté et sa position dans \vAttr{} des cas.
\end{itemize}
Les résultats par catégorie, les aires sous la courbe ROC et les taux de fausses alertes réalisés sont en
\cref{app-results}.

\subsection{Ce que deviennent les idées de départ}
\paragraph{B-spline et mélange gaussien.} L'idée de départ résumait chaque courbe par ses coefficients B-spline, par
région et catégorie, et regroupait ces coefficients par un mélange gaussien. Deux propriétés limitent cette approche.

\begin{proposition}[Ce que les coefficients ne voient pas]\label{prop-leverage}
Soit $Z\in\M_{N,n_s}(\R)$ la matrice des $n_s$ fonctions B-spline aux piliers et $H=Z(\tr{Z}Z)^{-1}\tr{Z}$ la
matrice de projection associée. L'ajustement par moindres carrés d'un pic unitaire $e_k$ a pour énergie
$\norm{He_k}^2=h_{kk}$, avec $\sum_kh_{kk}=n_s$. Une fraction $1-h_{kk}$ de l'énergie du pic reste dans le résidu et échappe à toute variable
calculée à partir des coefficients.
\end{proposition}
\begin{proof}
$H$ est symétrique et idempotente, donc $\norm{He_k}^2=\tr{e_k}He_k=h_{kk}$, et $\sum_kh_{kk}=\operatorname{tr}H=n_s$.
Les coefficients $(\tr{Z}Z)^{-1}\tr{Z}e_k$ déterminent $He_k$ et rien d'autre.
\end{proof}

Pour les six B-splines cubiques du banc d'essai et les onze piliers :
\begin{center}\input{generated/\lang/tab_leverage}\end{center}
Aux piliers intérieurs, les coefficients ne captent que \vLevMin{} à \vLevMax{} de l'énergie d'un pic, et la spline
répartit ce qu'elle capte sur plusieurs coefficients. Une anomalie sur un seul pilier, le motif le plus fréquent, est
diluée avant même que le regroupement commence.

\begin{proposition}[Pourquoi un mélange sur l'écart signale les erreurs FO]\label{prop-gmm-fo}
Prenons pour variables les coefficients $c=Z^+(q^{\tot}-q^{\fo})$, où $Z^+$ est la pseudo-inverse de $Z$. Un
décalage parallèle du repo FO de $u$ les déplace de $-u\,Z^+\one$, une direction fixe. Si le mélange est ajusté sur un
historique sain, où Totem et le FO concordent au bruit près, sa densité décroît dans cette direction : son taux
d'alerte sur les erreurs de marquage FO tend vers $1$ quand $|u|$ augmente, quelle que soit la qualité de la courbe
Totem.
\end{proposition}
\begin{proof}
$c$ est linéaire en $q^{\tot}-q^{\fo}$, qui se déplace de $-u\one$. Chaque composante gaussienne a une variance
finie dans la direction $Z^+\one$ : la densité du mélange, et avec elle le score de vraisemblance, tend vers zéro dans
cette direction.
\end{proof}
Le détecteur structurel n'a pas ce défaut, car $X\beta$ contient l'écart de niveau : il est retiré de chaque test et
publié à part. Le banc d'essai confirme les deux énoncés (\vFoGmmDiff{} contre \vFoStruct).

\paragraph{Petits réseaux de neurones par type d'anomalie.}
\begin{proposition}[Un réseau ne peut qu'apprendre le filtre adapté]\label{prop-nn}
Après blanchiment et projection, $\eh\sim\N(0, \Pi)$ sous $H_0$ et $\eh\sim\N(a\gt, \Pi)$ sous $H_1$, pour une forme $g$ et
une amplitude $a$ connues. Alors
\begin{equation}\label{eq-nn}
\ln\frac{\mathbb P(H_1\mid\eh)}{\mathbb P(H_0\mid\eh)}=a\sprod{\gt}{\eh}-\frac{a^2}{2}\norm{\gt}^2
+\ln\frac{\mathbb P(H_1)}{\mathbb P(H_0)},
\end{equation}
une fonction affine du filtre adapté $\sprod{\gt}{\eh}$. Un classifieur entraîné par entropie croisée sur assez de
données converge vers cette probabilité a posteriori ; avec une amplitude et un signe inconnus, on retrouve la
statistique $|z_g|$ du \cref{thm-glrt}.
\end{proposition}
\begin{proof}
Les deux lois gaussiennes ont la même covariance : les termes quadratiques de leur rapport se compensent. L'entropie
croisée est minimisée par la vraie probabilité a posteriori ; maximiser \eqref{eq-nn} en $a$ donne
$\sprod{\gt}{\eh}^2/(2\norm{\gt}^2)$.
\end{proof}

Pour les mécanismes que l'on sait nommer, un réseau ne peut donc au mieux qu'apprendre le filtre adapté, que le
détecteur structurel écrit directement, avec un taux de fausses alertes exact et une puissance connue.

\paragraph{Où ces outils restent utiles.}
\begin{itemize}
\item \textbf{Splines pénalisées dans le modèle nul}, si $\beta_0+\beta_1\ln T$ se révèle trop rigide sur données réelles.
\item \textbf{Mélanges gaussiens sur les sorties du détecteur} ($\hat\beta_0$, $\hat\beta_1$, les $z_g$, $\norm{\eh}^2$,
les rapports de dispersion), pour apprendre les régimes habituels par région et catégorie dans un espace de petite
dimension où chaque coordonnée a un sens.
\item \textbf{Région $\times$ catégorie} pour regrouper la calibration de $\kappa$ et de l'échelle de niveau.
\item \textbf{Réseaux de neurones}, une fois qu'existe un historique d'incidents étiquetés et des motifs que les formes
types ne couvrent pas ; ils doivent alors être entraînés sur les résidus $\eh$ et sur des injections synthétiques, pas
sur les courbes brutes.
\end{itemize}

\subsection{Limites}
\begin{attention}
Le marché synthétique reproduit les mécanismes par construction ; il ne montre pas que les données Totem réelles ne
contiennent que ceux-là. Les données réelles comporteront aussi des contributeurs figés, des opérations sur titres et
des dividendes exceptionnels, de l'exercice américain dans les cotations d'actions et des erreurs de saisie. L'échelle
$\kappa$ absorbe un niveau de bruit mal spécifié, pas une structure mal spécifiée, et les amplitudes injectées, bien que
plausibles, ne sont calibrées sur aucun historique. Le banc d'essai montre que la méthode fonctionne quand les
mécanismes sont présents et qu'elle est calibrée quand ils sont absents ; la validation sur douze à vingt-quatre mois
d'historique réel, avec des incidents étiquetés, est l'étape suivante.
\end{attention}

\begin{retenir}
\begin{itemize}[leftmargin=12pt]
\item À taux de fausses alertes égal, le détecteur structurel détecte autant que la meilleure alternative en moyenne,
et nettement plus sur les forwards aberrants.
\item Il ne prend pas une erreur de marquage FO pour une courbe aberrante, contrairement à un mélange gaussien sur
l'écart Totem--FO, et il explique ses alertes.
\item Les B-splines, les mélanges et les réseaux de neurones restent utiles, mais en aval du modèle structurel, pas sur
les courbes brutes.
\end{itemize}
\end{retenir}
